\subsection{Results}
\phantomsection\label{sec:ch02111}

Chapter~\ref{ch:02} may be read as a single logical chain. Impressed sources on bounded support,
together with charge continuity and declared boundary data, determine a harmonic Maxwell field
through $\mathcal{S}_\omega$. Outgoing Sommerfeld selection and radiative projection
$\mathcal{R}_\omega=\Pi_{\mathrm{rad}}\circ\mathcal{S}_\omega$ isolate the transverse content
that carries net export to infinity \cite{maxwell1873,stratton1941,sommerfeld1949}. Radiation is
therefore not an independent physical process added to Maxwell theory; it is the qualified
subset of solutions whose far-zone asymptotics support conserved power and pattern functionals.

Several structural results recur. The radiation map possesses nontrivial null spaces:
physically realizable impressed currents can exist whose radiative image vanishes while local
fields remain substantial \cite{devaney2004}. Electrical compactness and finite observation rank
compress accessible angular degrees of freedom through spectral decay distinct from nullity
\cite{chu1948,hansen1988}. Near-to-far reporting requires outgoing selection before asymptotic
extraction; reactive and evanescent admixtures can inflate tangential data on enclosing shells
without altering $a(\widehat{\mathbf{k}})$ \cite{stratton1941}. Reciprocity constrains bilinear
coupling on $\operatorname{ran}\mathcal{R}_\omega$ but does not remove finite-rank observation
limits or enlarge geometrically suppressed directions \cite{stratton1941,jackson1999}.

Principle~\ref{prin:ch2.radiative-accessibility} states that $\mathcal{F}_{\mathrm{rad}}$ must be
fixed before harmonic readouts; finite rank caps recoverable export \cite{stratton1941}. Results
proved throughout Sections~\ref{sec:ch20}--\ref{sec:ch210} presuppose Chapter~\ref{ch:01}
admissibility, outgoing-qualified radiation data, and the reactive-versus-radiative distinction of
Section~\ref{sec:ch01.2.4} \cite{stratton1941}. The completion condition
Eq.~\eqref{eq:ch2.continuity-admissibility} collects those requirements; report classes remain
those of Eq.~\eqref{eq:ch1.report-class-lock} \cite{stratton1941}.
\medskip
\begin{equation}
\label{eq:ch2.continuity-admissibility}
\mathcal{D}_{\mathrm{ch2}}=1
\iff
\left(
\mathcal{D}_{\mathrm{ch1}}=1\
\wedge\
\mathcal{R}_\omega[\Jimp]\ \text{is outgoing-qualified}\
\wedge\
\mathcal{D}_{J,\mathrm{reg}}=1
\right),
\end{equation}
with \(\mathcal{D}_{\mathrm{ch1}}\) from Eq.~\eqref{eq:ch1.continuity-admissibility} and
\(\mathcal{D}_{J,\mathrm{reg}}\) from Eq.~\eqref{eq:ch1.regime-qualified-decision}.
Let \(\mathcal{R}_{\mathrm{ch2}}\) denote theorems and propositions of
Sections~\ref{sec:ch20}--\ref{sec:ch210} proved under \(\mathcal{D}_{\mathrm{ch2}}=1\)
\cite{stratton1941}.